\documentclass[12pt]{article}
\usepackage[utf8]{inputenc}
\usepackage[T1]{fontenc}
\usepackage{hyperref}
\title{Muji Yi Language Resources for K-12 Curriculum Development}
\author{Compiled by Claude Sonnet 5.5 (claude-sonnet-5-5) for Andr\'as Kornai}
\date{}
\begin{document}
\maketitle

\section{Introduction}
The code ymj (Muji Yi) is retired. According to Glottolog, it was retired on 14 January 2008 and split into five languages after survey work in Yunnan in 2005--2006. The successors are Southern Muji [ymc], Moji [ymi], Qila Muji [ymq], Northern Muji [ymx] and Muzi [ymz]. No resources can be assigned to ymj itself. Resources must be sought under the successor codes, and each of those is a separate language. No educational materials were found for any of them.

\section{Primary Educational Resources}
None found for ymj or its successors.

\section{Linguistic and Reference Materials}
\begin{itemize}
\item Glottolog, ymj record (retired code): \url{https://glottolog.org/resource/languoid/iso/ymj}
\item Wikipedia, Moji language (ymi): \url{https://en.wikipedia.org/wiki/Moji_language}
\item Wikipedia, Muzi language (ymz): \url{https://en.wikipedia.org/wiki/Muzi_language}
\end{itemize}

\section{Religious and Community Resources}
\begin{itemize}
\item Joshua Project, Northern Muji: \url{https://joshuaproject.net/people_groups/20606}
\item Joshua Project, Southern Muji: \url{https://joshuaproject.net/people_groups/18613/CH}
\item Asia Harvest Muji people-group profile (PDF): \url{https://people-groups.asiaharvest.org/people-group-profiles/yunnan/muji.pdf}
\end{itemize}

\section{Digital and Audio Resources}
None found.

\section{Curriculum Development Context}
Because the code is retired, any curriculum work should target a specific successor language, such as Southern Muji (the most widely spoken, with about 29,000 speakers according to one source) or Northern Muji. The Muji are officially counted within the Yi nationality.

\end{document}